
\documentclass{article}
\usepackage{colortbl}
\usepackage{makecell}
\usepackage{multirow}
\usepackage{supertabular}

\begin{document}

\newcounter{utterance}

\centering \large Interaction Transcript for game `eqbench', experiment `eq\_bench', episode 36 with reflexion{-}llama{-}3.1{-}8b{-}instruct.
\vspace{24pt}

{ \footnotesize  \setcounter{utterance}{1}
\setlength{\tabcolsep}{0pt}
\begin{supertabular}{c@{$\;$}|p{.15\linewidth}@{}p{.15\linewidth}p{.15\linewidth}p{.15\linewidth}p{.15\linewidth}p{.15\linewidth}}
    \# & \multicolumn{2}{c}{Player} && \multicolumn{2}{c}{Game Master} \\
    \hline

    \theutterance \stepcounter{utterance}  
    & & & \multicolumn{4}{p{0.6\linewidth}}{
        \cellcolor[rgb]{0.9,0.9,0.9}{
            \makecell[{{p{\linewidth}}}]{
                \texttt{\tiny{[P1$\langle$GM]}}
                \texttt{You have attempted this same task before and did not succeed. These are the plans you wrote after those attempts.} \\
\texttt{Trial \#1: New plan: } \\
\\ 
\texttt{Upon re{-}examining the dialogue, I realized that I focused too much on the surface{-}level conversation and didn't consider the underlying emotional dynamics. I should have paid closer attention to Deirdre's emotional state and Scott's responses to her concerns. } \\
\\ 
\texttt{In my next attempt, I will analyze the dialogue more deeply, considering the following:} \\
\\ 
\texttt{1. Deirdre's emotional state: I will identify her primary emotions and concerns, such as feeling abandoned, rejected, and worried about her future.} \\
\texttt{2. Scott's responses: I will evaluate how Scott's words and actions affect Deirdre's emotions, such as his dismissive tone, his emphasis on his own needs, and his lack of empathy.} \\
\texttt{3. Power dynamics: I will consider the power imbalance between Deirdre and Scott, with Scott being the one leaving and Deirdre being the one left behind.} \\
\\ 
\texttt{With this new understanding, I will reassess the emotions Deirdre is likely to feel at the end of the dialogue.} \\
\\ 
\\ 
\\ 
\texttt{Your task is to predict the likely emotional responses of a character in this dialogue:} \\
\\ 
\texttt{Deirdre: Scott, have you thought about this? Really thought about it?} \\
\texttt{Scott: Yes, mom. I have. } \\
\texttt{Deirdre: You're just going to leave me here, all alone?} \\
\texttt{Scott: That's not fair, mom. You know it's not like that.} \\
\texttt{Deirdre: Isn't it? Your father left me, and now you're doing the same.} \\
\texttt{Scott: Dad didn't leave, mom. He died. There's a difference.} \\
\texttt{Deirdre: But the end result's the same, isn't it? I'm left alone.} \\
\texttt{Scott: I'm not abandoning you. I'm just trying to live my life.} \\
\texttt{Deirdre: And what about my life, Scott? Who's going to take care of me?} \\
\texttt{Scott: You've got Mrs. Henderson next door, and I'll hire a nurse if need be.} \\
\texttt{Deirdre: I don't want a nurse, Scott. I want my son.} \\
\texttt{Scott: I can't be your crutch forever, mom.} \\
\texttt{Deirdre: Is that what I am to you now? A burden?} \\
\texttt{Scott: No, mom. But I can't keep sacrificing my dreams for your fears.} \\
\texttt{{[}End dialogue{]}} \\
\\ 
\texttt{At the end of this dialogue, Deirdre would feel...} \\
\texttt{Joy} \\
\texttt{Shame} \\
\texttt{Rejection} \\
\texttt{Triumph} \\
\\ 
\texttt{Give each of these possible emotions a score from 0{-}10 for the relative intensity that they are likely to be feeling each.} \\
\\ 
\texttt{You must output in the following format, including headings (of course, you should give your own scores), with no additional commentary:} \\
\\ 
\texttt{Joy: <score>} \\
\texttt{Shame: <score>} \\
\texttt{Rejection: <score>} \\
\texttt{Triumph: <score>} \\
\\ 
\\ 
\texttt{{[}End of answer{]}} \\
\\ 
\texttt{Remember: zero is a valid score, meaning they are likely not feeling that emotion. You must score at least one emotion > 0.} \\
\\ 
\texttt{Your answer:} \\
            }
        }
    }
     \\ \\

    \theutterance \stepcounter{utterance}  
    & \multicolumn{4}{p{0.6\linewidth}}{
        \cellcolor[rgb]{0.561,0.737,0.561}{
            \makecell[{{p{\linewidth}}}]{
                \texttt{\tiny{[P1$\rangle$GM]}}
                \texttt{Joy: 0} \\
\texttt{Shame: 2} \\
\texttt{Rejection: 9} \\
\texttt{Triumph: 0} \\
            }
        }
    }
    & & \\ \\

    \theutterance \stepcounter{utterance}  
    & & & \multicolumn{2}{p{0.3\linewidth}}{
        \cellcolor[rgb]{0.9,0.9,0.9}{
            \makecell[{{p{\linewidth}}}]{
                \texttt{\tiny{[GM$|$GM]}}
                \texttt{Joy: 0} \\
\texttt{Shame: 2} \\
\texttt{Rejection: 9} \\
\texttt{Triumph: 0} \\
            }
        }
    }
    & & \\ \\

    \theutterance \stepcounter{utterance}  
    & & & \multicolumn{2}{p{0.3\linewidth}}{
        \cellcolor[rgb]{0.9,0.9,0.9}{
            \makecell[{{p{\linewidth}}}]{
                \texttt{\tiny{[GM$|$GM]}}
                \texttt{Joy: 0} \\
\texttt{Shame: 4} \\
\texttt{Rejection: 8} \\
\texttt{Triumph: 0} \\
            }
        }
    }
    & & \\ \\

    \theutterance \stepcounter{utterance}  
    & & & \multicolumn{2}{p{0.3\linewidth}}{
        \cellcolor[rgb]{0.9,0.9,0.9}{
            \makecell[{{p{\linewidth}}}]{
                \texttt{\tiny{[GM$|$GM]}}
                \texttt{game\_result = LOSE} \\
            }
        }
    }
    & & \\ \\

    \theutterance \stepcounter{utterance}  
    & & \multicolumn{4}{p{0.6\linewidth}}{
        \cellcolor[rgb]{0.95,0.95,0.95}{
            \makecell[{{p{\linewidth}}}]{
                \texttt{\tiny{[P1$\leftrightarrow$P]}}
                \texttt{Post{-}episode Reflexion (not sent to GM):} \\
\texttt{New plan: } \\
\\ 
\texttt{Upon re{-}examining the dialogue, I realized that I focused too much on the surface{-}level conversation and didn't consider the underlying emotional dynamics. I should have paid closer attention to Deirdre's emotional state and Scott's responses to her concerns. } \\
\\ 
\texttt{In my next attempt, I will analyze the dialogue more deeply, considering the following:} \\
\\ 
\texttt{1. Deirdre's emotional state: I will identify her primary emotions and concerns, such as feeling abandoned, rejected, and worried about her future.} \\
\texttt{2. Scott's responses: I will evaluate how Scott's words and actions affect Deirdre's emotions, such as his dismissive tone, his emphasis on his own needs, and his lack of empathy.} \\
\texttt{3. Power dynamics: I will consider the power imbalance between Deirdre and Scott, with Scott being the one leaving and Deirdre being the one left behind.} \\
\\ 
\texttt{With this new understanding, I will reassess the emotions Deirdre is likely to feel at the end of the dialogue. I will also consider the following specific actions to improve my analysis:} \\
\\ 
\texttt{{-} I will pay closer attention to Deirdre's language and tone, looking for emotional cues such as phrases that indicate fear, anxiety, or sadness.} \\
\texttt{{-} I will evaluate Scott's responses more critically, considering how his words and actions may be perceived as insensitive or dismissive.} \\
\texttt{{-} I will consider the long{-}term implications of Scott's decision to leave, and how it may affect Deirdre's emotional state in the future.} \\
\\ 
\texttt{With this new approach, I will reassess the emotions Deirdre is likely to feel at the end of the dialogue, and provide a revised answer.} \\
            }
        }
    }
    & \\ \\

\end{supertabular}
}

\end{document}
